

% ======================================================================
\Needspace{12\baselineskip}
\section{Training Loop}
% ======================================================================

\problem{data\_loading}{Implement data loading}{2}

\deliverable{Write a function that takes a numpy array, batch size, context length, and device, and returns a pair of tensors.}

% Implementation summary:
Sample $B$ starting indices independently
and uniformly with replacement. For $N$ dataset tokens and context length $T$,
\[
s_b\sim\operatorname{Uniform}\{0,\ldots,N-T-1\},\qquad
X_{b,t}=x_{s_b+t},\quad Y_{b,t}=x_{s_b+t+1},\quad 0\leq t<T.
\]
Thus $X$ and $Y$ have shape $(B,T)$, with targets shifted one token forward.
\par\noindent
\begin{minipage}{\linewidth}
\begin{lstlisting}[caption={Sampling input and target windows}]
starts = np.random.randint(0, len(dataset) - context_length, size=batch_size)
inputs_np = np.stack([dataset[s:s + context_length] for s in starts])
targets_np = np.stack([dataset[s + 1:s + context_length + 1] for s in starts])
inputs = torch.tensor(inputs_np, dtype=torch.long, device=device)
targets = torch.tensor(targets_np, dtype=torch.long, device=device)
\end{lstlisting}
\end{minipage}\par
Only sampled windows are materialized; the dataset can remain memory-mapped.


\problem{checkpointing}{Implement model checkpointing}{1}

\deliverable{Implement \texttt{save\_checkpoint} and \texttt{load\_checkpoint} functions.}

% Implementation summary:
A \texttt{torch.save} dictionary contains
the following state:
\begin{center}
\begin{tabular}{ll}
\toprule
Field & Purpose \\
\midrule
\texttt{model}, \texttt{optimizer} & Restore weights and AdamW state \\
\texttt{iteration} & Resume at the next unfinished update \\
\texttt{model\_config} (optional) & Reconstruct the model for generation \\
\texttt{numpy\_rng\_state}, \texttt{torch\_rng\_state} & Restore CPU random generators \\
\texttt{mps\_rng\_state}, \texttt{cuda\_rng\_states} & Restore available accelerator generators \\
\bottomrule
\end{tabular}
\end{center}
Loading restores model, optimizer, and random generator states and returns
the iteration count. The training script always supplies the model
configuration, enabling inference to reconstruct the architecture:
\par\noindent
\begin{minipage}{\linewidth}
\begin{lstlisting}[caption={Reconstructing a model from a training checkpoint}]
checkpoint = torch.load(path, map_location="cpu", weights_only=False)
model = TransformerLM(**checkpoint["model_config"])
model.load_state_dict(checkpoint["model"])
\end{lstlisting}
\end{minipage}\par


\problem{training\_together}{Put it together}{4}

\deliverable{Write a script that runs a training loop to train your model on user-provided input.}

% Implementation summary:
The training script exposes data paths,
model and optimizer hyperparameters, device, random seed, logging intervals,
and checkpoint paths through command-line arguments. Training and validation
token files are opened as read-only \texttt{uint16} NumPy memory maps to
avoid loading the entire corpora into memory.

Each update sets the warmup/cosine learning rate before computing the loss.
The core update, after setting the learning rate, is:
\par\noindent
\begin{minipage}{\linewidth}
\begin{lstlisting}[caption={Training update}]
optimizer.zero_grad()
inputs, targets = get_batch(
    train_dataset, args.batch_size, args.context_length, args.device
)
logits = model(inputs)
loss = cross_entropy(
    logits.reshape(-1, logits.shape[-1]), targets.reshape(-1)
)
loss.backward()
gradient_clipping(model.parameters(), args.max_l2_norm)
optimizer.step()
\end{lstlisting}
\end{minipage}\par
Validation averages losses over sampled batches under \texttt{eval()} with
gradients disabled, then restores the previous mode. Configuration, training
loss, validation loss, and learning rate are logged to the console and
\texttt{train.log}. Checkpoints are saved periodically and at the end;
resuming restores saved states using the supplied training configuration.


